\section{The source's finite size-principle calculation}
\label{sec:p11-annex-rationality-12-3}

\begin{annexitem}{annex:p11.r123.size-principle}{thm:p11-r123-size-principle}{Fair-die likelihoods and posterior odds}
\begin{proof}
Under \(H_d\), an admissible face has probability \(1/d\) and
an inadmissible face has probability zero. Independence of the retained
die's rolls makes the ordered-sequence likelihood the product of those
probabilities. Thus it is \(d^{-n}\) when all faces are admissible,
and zero otherwise. The empty product is one.

Prior selection is uniform over the four hypotheses. Bayes' rule gives
\[
 p_d(x)=\frac{(1/4)L_d(x)}{\sum_{e\in D}(1/4)L_e(x)}
        =\frac{L_d(x)}{\sum_{e\in D}L_e(x)}
\]
whenever the observation has positive marginal probability. Every face
in the stated ten-roll sequence lies between 1 and 4, so every candidate
is compatible. Substituting \(n=10\) proves the displayed formula.
With common denominator \(601368567754849\), its four posterior
numerators in increasing die-size order are
\[
590490000000000,\quad10240000000000,\quad
576650390625,\quad61917364224.
\]
They sum to the denominator. Division gives the reported approximations;
the four-sided hypothesis has the largest posterior.

For two compatible candidates \(d<e\), division of their posterior
probabilities cancels the same positive normalizing sum and gives
\(p_d(x)/p_e(x)=d^{-n}/e^{-n}=(e/d)^n\).
For \(n>0\) this exceeds one; for \(n=0\), every likelihood is
one and the posterior is the uniform prior. If a candidate cannot
generate one observation, its zero likelihood makes its posterior zero
provided some other candidate can generate the data. If no candidate
can generate the entire sequence, the marginal probability is zero and
conditioning by this formula is undefined: no arbitrary posterior is
assigned. Unequal prior odds would multiply, not disappear from, the
posterior-odds comparison. The result therefore relies on its declared
sampling law and priors.
\end{proof}
\end{annexitem}
